\documentclass[12pt, a4paper]{report}

% ===================================================================
% PAQUETES ESENCIALES
% ===================================================================

% Codificación y lenguaje
\usepackage[utf8]{inputenc}
\usepackage[spanish]{babel}

% Matemáticas
\usepackage{amsmath}
\usepackage{amssymb}
\usepackage{amsfonts}

% Gráficos e imágenes
\usepackage{graphicx}
\graphicspath{{../assets/images/}}

% Configuración de página
\usepackage{geometry}
\geometry{a4paper, margin=1in}

% Enlaces e hipervínculos
\usepackage{hyperref}
\hypersetup{
    colorlinks=true,
    linkcolor=blue,
    filecolor=magenta,
    urlcolor=cyan,
    pdftitle={Título del Reporte},
    pdfauthor={Autor},
}

% Tablas mejoradas
\usepackage{booktabs}
\usepackage{longtable}

% Código fuente (opcional)
\usepackage{listings}
\usepackage{xcolor}

\lstset{
    basicstyle=\ttfamily\small,
    breaklines=true,
    frame=single,
    numbers=left,
    numberstyle=\tiny,
    backgroundcolor=\color{gray!10},
}

% ===================================================================
% INFORMACIÓN DEL DOCUMENTO
% ===================================================================

\title{
    \textbf{Título del Reporte Técnico} \\
    \large Subtítulo o descripción adicional
}
\author{
    Nombre del Autor \\
    Organización/Empresa \\
    \texttt{email@ejemplo.com}
}
\date{\today}

% ===================================================================
% INICIO DEL DOCUMENTO
% ===================================================================

\begin{document}

\maketitle

\begin{abstract}
Resumen ejecutivo del reporte. Describe brevemente el propósito, alcance y conclusiones principales del documento.
\end{abstract}

\tableofcontents
\listoffigures
\listoftables

\chapter{Introducción}

\section{Contexto}
Describe el contexto del reporte.

\section{Objetivos}
Lista los objetivos del reporte.

\section{Alcance}
Define el alcance del documento.

\chapter{Marco Teórico}

Presenta los conceptos teóricos relevantes.

\chapter{Metodología}

Describe la metodología utilizada.

\chapter{Resultados}

\section{Resultados principales}
Presenta los resultados obtenidos.

\section{Análisis de resultados}
Analiza los resultados.

\chapter{Conclusiones y Recomendaciones}

\section{Conclusiones}
Resume las conclusiones principales.

\section{Recomendaciones}
Proporciona recomendaciones basadas en los resultados.

\appendix

\chapter{Anexos}

Material complementario, código fuente, tablas adicionales, etc.

\begin{thebibliography}{9}
\bibitem{ref1}
Referencia bibliográfica 1.

\bibitem{ref2}
Referencia bibliográfica 2.
\end{thebibliography}

\end{document}
